\chapter{Introduzione}

\begin{section}{Dispositivi Indossabili}
    Pur essendo presenti in campo medico da tempo, i dispositivi indossabili e impiantabili hanno subito una rapida evoluzione. Tra le possibili applicazioni si hanno diagnosi attraverso la registrazione di segnali e marcatori biologici, ma anche terapia attraverso stimoli elettrici o rilascio di medicinali. In particolare, la portabilit\`a abilita al funzionamento lontano dagli ospedali, rispondendo ai requisiti della telemedicina. Questo li rende particolarmente adatti al trattamento di condizioni croniche, sgravando sia strutture sanitarie, che pazienti. Permettono inoltre misure costanti, a lungo termine e non invasive. 

    I dispositivi indossabili, oltre a sensori e attuatori specifici per l'applicazione, richiedono anche componenti a bassa potenza come processori e antenne. Questi componenti sono necessari per svolgere funzioni di controllo automatico, elaborazione e trasmissione dei segnali misurati. Quando un modulo wireless \'e disponibile, invece di conservare i dati, \'e possibile inviarli a una hub intermedia o direttamente al destinatario finale in una struttura medica.

    La forma ideale \'e piccola abbastanza da svolgere le funzioni volute senza intralciare il paziente. Nelle installazioni invasive \'e essenziale che le dimensioni siano ridotte per abbassare i rischi e la complessit\'a dell'operazione. Anche nel caso non invasivo per\'o, si preferiscono dispositivi sottili e di piccole dimensioni che possono sostituire o integrarsi a capi di abbigliamento. L'aspetto non intrusivo o piacevole pu\'o incentivare l'uso regolare. Evitare dimenticanze \'e infatti particolarmente importante in pazienti a rischio che richiedono controllo costante. Avere numerosi punti di possibile applicazione permette di usare pi\'u dispositivi in combinazione.

    \begin{figure}[H]
        \includegraphics[width=1\textwidth]{indossabili.png}
        \centering
        \caption{a) Alcuni esempi di dispositivo medico indossabile. b) Flusso informativo dei dispositivi medici indossabili.}
        \label{fig:indossabili}
    \end{figure}
    
    Si vuole quindi approfondire uno degli aspetti pi\'u critici dei dispositivi medici indossabili e impiantabili, ovvero la loro alimentazione. Ogni tipologia avr\'a necessit\'a energetiche diverse da pochi microwatt a milliwatt, ma al momento hanno tutti in comune l'uso di batterie come fonte. Gli svantaggi di un sistema a batteria sono evidenti nell'autonomia limitata e il bisogno di ricarica o sostituzione, specialmente quando si parla di dispositivi impiantabili. Inoltre, i materiali usati nelle batterie sono pesanti e potenzialmente tossici, il che riduce la compatibilit\'a all'uso prolungato.
\end{section}

\begin{section}{Harvesting Energetico}
    Le tecniche di harvesting energetico consistono nel recupero di energia da fonti esterne al dispositivo. I vantaggi di queste tecnologie sono molteplici, ad esempio, estensione dell'autonomia in modo indefinito, riduzione di dimensioni e peso, flessibilit\'a e biocompatibilit\'a. Rappresentano una potenziale alternativa all'uso di batterie su dispositivi a bassa potenza.

    Le sorgenti principali a cui un harvester pu\'o attingere sono classificabili secondo lo schema \ref{fig:classificazione}. La potenza in ingresso a un harvester di dimensioni ridotte \'e necessariamente piccola, quindi quella prodotta sar\'a limitata anche se si lavora ad alto rendimento. Combinazioni di pi\'u sistemi, detti harvester ibridi, possono migliorare le prestazioni al costo di maggiore complessit\'a. La localizzazione del dispositivo sul corpo determina quale fonte energetica \'e pi\'u accessibile e quindi in gran parte anche la tecnologia di harvesting migliore.
    \begin{figure}[H]
        \includegraphics[width=0.9\textwidth]{classificazioni_fonti.png}
        \centering
        \caption{Classificazione delle fonti energetiche a cui pu\'o essere applicato un harvester per dispositivi indossabili.}
        \label{fig:classificazione}
    \end{figure}
    
    La maggior parte delle fonti in figura \ref{fig:classificazione} non \'e sempre disponibile e nemmeno prevedibile. \'E possibile per\'o usare harvester accoppiati a componenti di accumulo come batterie o supercondensatori quando \'e richiesta un'alimentazione regolare. Un'alternativa per monitorare segnali a rischio pi\'u basso \'e accendere il dispositivo solo quando l'energia \'e disponibile. Similmente, per aumentare l'efficienza, si possono alimentare i singoli moduli quando diventano necessari. La gestione intelligente dell'energia richiede una power managment unit.
\end{section}

\begin{section}{Appllicazioni Harvester Energetici}
    Esistono versioni indossabili per la misura non invasiva di molti parametri di interesse rilevante come ECG, EEG, pressione sanguigna, saturazione di ossigeno nel sangue, temperatura e frequenza respiratoria. Altri dispositivi sono in grado di raccogliere in maniera pi\'u o meno invasiva fluidi corporei come sangue, sudore o liquido interstiziale per intercettare problemi metabolici. Oltre a raccogliere dati, questi apparati medici possono anche partecipare attivamente alla cura, come nel caso di dosatori di medicinali, elettrostimolatori neuronali o muscolari, pacemaker e defibrillatori impiantabili.
    
    Monitorare facilmente e con costanza questi parametri consente di individuare lo sviluppo di alcune malattie in anticipo rispetto ai metodi tradizionali, aumentando le probabilit\'a di successo nel trattamento. Nonostante questo, nessun dispositivo medico indossabile commercializzato fa uso di energy harvesting. La ricerca recente ha per\'o dimostrato con esempi pratici l'efficacia di questa metodologia. Per dispositivi a uso esterno facilmente testabili, ma anche per impianti in vivo su animale. Alcuni esempi recenti sono: \begin{itemize}
        \item Un dispositivo in grado di misurare glicemia dal sudore, temperatura e HRV \'e stato sviluppato in \cite{mirlouContinuousGlycemicMonitoring2024}. Un harvester RF ne permette il funzionamento continuativo, supportando la batteria.
        \item Per trattare la miopia con esposizione graduale ai farmaci, in \cite{jiangSelfgeneratedElectricitydrivenDrug2024} \'e stata sviluppato un dosatore a forma di lente a contatto. La funzione di rilascio \'e automatizzata da un nanogeneratore piezoelettrico attivato dal battito delle palpebre.
        \item Un modulo autosufficiente che monitora respiri e temperatura \'e stato integrato in mascherine sanitarie in \cite{lanHighefficientIntelligentAntibacterial2024}. In \cite{simInstantDisinfectingFace}, un harvester triboelettrico \'e usato per creare un campo elettrico forte abbastanza da disinfettare le maschere.
        \item In \cite{anVivoFlexibleEnergy2024} viene dimostrato l'uso di un harvester piezoelettrico impiantato su tessuto cardiaco porcino.
    \end{itemize}
 
    Soluzioni con harvester hanno alcuni svantaggi rispetto a sola batteria. Tra questi, la bassa densit\'a energetica e l'affidabilit\'a incerta li rendono un opzione immatura per scopi medici critici. Questi sono gravosi soprattutto per dispositivi impiantabili che richiedono dimensioni minime e la certezza del funzionamento a lungo termine. I costi di produzione sono alti, dovuti anche alla mancanza di una catena manifatturiera paragonabile a prodotti altamente diffusi come le batterie.

    \'E giusto menzionare che \'e possibile applicare l'idea di energy harvester anche a dispositivi non mobili. In questo caso, naturalmente, non si pu\'o pi\'u attingere alle fonti energetiche del corpo umano, ma se ne hanno altre come vento o vibrazioni. Un'applicazione di interesse \'e usarli per recuperare l'energia di scarto da sistemi pi\'u grandi come macchinari industriali o mezzi di trasporto. In zone dove \'e necessario fare monitoraggi, ma manca distribuzione elettrica, possono essere una soluzione. Questo permetterebbe di azionare una rete distribuita di sensori o attuatori autonomi. Harvester impiegati in questo modo hanno meno restrizioni su dimensioni e affidabilit\'a. Alcuni esempi recenti sono: \begin{itemize}
        \item Monitoraggio della sicurezza e stato dei mezzi di trasporto come treni o auto trasformando l'energia delle vibrazioni \cite{liSmartRailwayTransportation, liuCompactHybridizedTriboelectricelectromagnetic2024}
        \item Una rete di generatori galleggianti che raccolgono l'energia meccanica delle onde per deionizzare l'acqua marina in loco \cite{renWavepoweredCapacitiveDeionization2024}. 
        \item Sensori per verificare lo stato di macchinari industriali, sfruttandone le vibrazioni \cite{alvarezruedaVibrationEnergyHarvesting2024, gaoHybridGeneratorEfficient2024}, o il calore \cite{deoliveiraDevelopmentHybridEnergy2024}.
    \end{itemize}
    
\end{section}
